\chapter{MÓDULO DE GESTIÓN DE DESAFÍOS}

Con este módulo el organizador prepara cada edición del Desafío Bebras. Tiene dos partes que dependen una de la otra: el banco de tareas, donde se redactan las preguntas con su respuesta y su explicación, y los desafíos, que toman preguntas del banco, las ordenan por categoría y definen cuándo y cómo se rinden. Incluye además la práctica, donde cualquier estudiante se ejercita con tareas liberadas.

Las tareas Bebras no son solo de opción múltiple: el estudiante también arrastra piezas sobre un dibujo, señala una zona de una imagen, elige el estado de cada casilla de una cuadrícula o completa los huecos de un texto. Por eso el módulo ofrece seis tipos de respuesta, cada uno con su editor y su forma de corrección.

\section{PLANIFICACIÓN DEL MÓDULO}

Esta sección fija el propósito del módulo, nombra sus historias de usuario y reparte su entrega entre las Iteraciones 3 y 4.

\subsection{Objetivo del módulo}

Permitir la creación, categorización y administración de tareas de pensamiento computacional según rangos de edad y niveles de dificultad, y organizar con ellas desafíos para varias categorías a la vez.

\subsection{Historias de usuario}

Las historias responden a las necesidades del administrador, que prepara las tareas y los desafíos; del estudiante y del maestro, que practican; y del visitante, que consulta los desafíos vigentes. Cada una representa una sola necesidad y lleva criterios de aceptación que permiten decidir cuándo está cumplida \parencite{cohn2004}; se agrupan según la actividad del usuario, siguiendo el mapeo de historias de \textcite{patton2014}. La prioridad 1 marca lo imprescindible para realizar un desafío y la 2, las mejoras que no lo impiden; el esfuerzo se mide en puntos, con la escala del módulo de gestión de contenido. La Tabla~\ref{tab:historias_desafios} las lista; las tarjetas completas están en el Anexo A.

{\estilotabla
\begin{longtable}{l >{\raggedright\arraybackslash}p{6.9cm} >{\raggedright\arraybackslash}p{2.4cm} c c}
\arrayrulecolor{huazul}\specialrule{1.4pt}{0pt}{0pt}
\rowcolor{hufondo}\textbf{ID} & \textbf{Historia} & \textbf{Usuario} & \textbf{Prioridad} & \textbf{Esfuerzo} \\
\arrayrulecolor{hulinea}\midrule
\endhead
HU-18 & Registrar una tarea de opción múltiple & Administrador & 1 & 13 \\
HU-19 & Registrar una tarea de respuesta corta & Administrador & 1 & 3 \\
HU-20 & Registrar una tarea de arrastrar y soltar & Administrador & 1 & 13 \\
HU-21 & Registrar una tarea de zonas sobre una imagen & Administrador & 2 & 13 \\
HU-22 & Registrar una tarea de completar huecos & Administrador & 2 & 13 \\
HU-23 & Registrar una tarea de estados por casilla & Administrador & 2 & 8 \\
HU-24 & Explicar la solución & Administrador & 1 & 5 \\
HU-25 & Clasificar una tarea & Administrador & 1 & 3 \\
HU-26 & Probar una tarea antes de usarla & Administrador & 1 & 5 \\
HU-27 & Decidir quién puede ver una tarea & Administrador & 1 & 3 \\
HU-28 & Encontrar una tarea en el banco & Administrador & 2 & 3 \\
HU-29 & Borrar una tarea que ya no sirve & Administrador & 2 & 3 \\
HU-30 & No alterar una tarea mientras se rinde & Administrador & 1 & 3 \\
HU-31 & Crear un desafío & Administrador & 1 & 5 \\
HU-32 & Elegir las preguntas de cada categoría & Administrador & 1 & 8 \\
HU-33 & Programar las fechas del desafío & Administrador & 1 & 8 \\
HU-34 & Ajustar el puntaje & Administrador & 2 & 3 \\
HU-35 & Decidir qué ve el estudiante al terminar & Administrador & 2 & 5 \\
HU-36 & Mezclar el orden de las opciones & Administrador & 2 & 3 \\
HU-37 & Publicar el desafío & Administrador & 1 & 5 \\
HU-38 & Recorrer el desafío como un estudiante & Administrador & 2 & 5 \\
HU-39 & Pausar un desafío en curso & Administrador & 2 & 8 \\
HU-40 & Practicar sin crear una cuenta & Estudiante & 1 & 8 \\
HU-41 & Consultar el avance de la práctica & Estudiante & 2 & 3 \\
HU-42 & Armar una práctica para un grupo & Maestro & 2 & 8 \\
HU-43 & Rendir la práctica de un maestro & Estudiante & 2 & 3 \\
HU-44 & Ver qué desafíos hay & Visitante & 1 & 5 \\
HU-45 & Ir directo a la acción de cada desafío & Visitante & 2 & 3 \\
\arrayrulecolor{hulinea}\midrule
\multicolumn{4}{r}{\textbf{Total de puntos}} & \textbf{168} \\
\arrayrulecolor{hulinea}\bottomrule
\rowcolor{white}\caption{Historias de usuario del módulo de gestión de desafíos}\label{tab:historias_desafios}\\
\end{longtable}
\arrayrulecolor{black}
}

\subsection{Plan de entregas y alcance del módulo}

La primera iteración construye el banco de tareas, porque un desafío no puede armarse sin preguntas; la segunda organiza los desafíos, la práctica y su presentación al público. La Tabla~\ref{tab:iteraciones_desafios} muestra el alcance y las historias de cada una.

\begin{table}[H]
\estilotabla
\begin{tabularx}{\linewidth}{l >{\raggedright\arraybackslash}X >{\raggedright\arraybackslash}p{3.6cm}}
\arrayrulecolor{huazul}\specialrule{1.4pt}{0pt}{0pt}
\rowcolor{hufondo}\textbf{Iteración} & \textbf{Alcance} & \textbf{Historias} \\
\arrayrulecolor{hulinea}\midrule
Iteración 3 (6 al 17 de julio) & Banco de tareas: redacción y explicación, seis tipos de respuesta, dificultad por categoría, probador, visibilidad y búsqueda & \mbox{HU-18} a \mbox{HU-30} \\
Iteración 4 (20 al 31 de julio) & Desafíos por categoría con sus preguntas, fechas, puntaje, publicación y pausa; práctica; portada pública & \mbox{HU-31} a \mbox{HU-45} \\
\arrayrulecolor{hulinea}\bottomrule
\end{tabularx}
\arrayrulecolor{black}
\caption[Plan de entregas del módulo de gestión de desafíos]{Distribución de historias del módulo de gestión de desafíos por iteración}
\label{tab:iteraciones_desafios}
\end{table}

\section{ITERACIÓN 3}

La Iteración 3 construyó el banco de tareas. Se describen su planificación, el diseño de la tarea y de sus tipos de respuesta, la codificación, las pruebas y los resultados.

\subsection{Planificación de la iteración}

La iteración debía permitir que el organizador redactara una tarea completa, eligiera cómo la responde el estudiante y la probara antes de usarla, de modo que las 43 tareas del cuadernillo Bebras 2024 pudieran representarse con fidelidad. Arrastrar y soltar, zonas sobre una imagen y completar huecos concentraron el mayor esfuerzo, por su interacción directa sobre imágenes y textos.

La Tabla~\ref{tab:hus_iteracion_3} lista las historias de la iteración; sus tarjetas, con los criterios de aceptación, están en el Anexo A.

\begin{table}[H]
\estilotabla
\begin{tabularx}{\linewidth}{l >{\raggedright\arraybackslash}X >{\raggedright\arraybackslash}p{2.8cm} c c}
\arrayrulecolor{huazul}\specialrule{1.4pt}{0pt}{0pt}
\rowcolor{hufondo}\textbf{ID} & \textbf{Historia} & \textbf{Usuario} & \textbf{Prioridad} & \textbf{Esfuerzo} \\
\arrayrulecolor{hulinea}\midrule
HU-18 & Registrar una tarea de opción múltiple & Administrador & 1 & 13 \\
HU-19 & Registrar una tarea de respuesta corta & Administrador & 1 & 3 \\
HU-20 & Registrar una tarea de arrastrar y soltar & Administrador & 1 & 13 \\
HU-21 & Registrar una tarea de zonas sobre una imagen & Administrador & 2 & 13 \\
HU-22 & Registrar una tarea de completar huecos & Administrador & 2 & 13 \\
HU-23 & Registrar una tarea de estados por casilla & Administrador & 2 & 8 \\
HU-24 & Explicar la solución & Administrador & 1 & 5 \\
HU-25 & Clasificar una tarea & Administrador & 1 & 3 \\
HU-26 & Probar una tarea antes de usarla & Administrador & 1 & 5 \\
HU-27 & Decidir quién puede ver una tarea & Administrador & 1 & 3 \\
HU-28 & Encontrar una tarea en el banco & Administrador & 2 & 3 \\
HU-29 & Borrar una tarea que ya no sirve & Administrador & 2 & 3 \\
HU-30 & No alterar una tarea mientras se rinde & Administrador & 1 & 3 \\
\arrayrulecolor{hulinea}\midrule
\rowcolor{white}\multicolumn{4}{r}{\textbf{Total de puntos}} & \textbf{88} \\
\arrayrulecolor{hulinea}\bottomrule
\end{tabularx}
\arrayrulecolor{black}
\caption[Historias de la Iteración 3]{Historias de la Iteración 3}
\label{tab:hus_iteracion_3}
\end{table}


\subsection{Diseño}

Una tarea tiene un enunciado, una pregunta y una explicación, armados con bloques de texto con formato e imágenes, y una dificultad para cada categoría en la que se puede usar. Su respuesta separa lo que el estudiante necesita para responder de lo que el servidor usa para corregir: los tipos más recientes (zonas, estados por casilla y huecos) lo guardan en una configuración pública y una clave privada, y los primeros, en campos propios de cada tipo. El servidor arma para el estudiante una versión de la tarea que nunca incluye la solución. La Tabla~\ref{tab:tipos_respuesta} resume cómo responde el estudiante y cómo se corrige cada uno de los seis tipos.

\begin{table}[H]
\estilotabla
\begin{tabularx}{\linewidth}{>{\raggedright\arraybackslash}p{3.3cm} >{\raggedright\arraybackslash}X >{\raggedright\arraybackslash}X}
\arrayrulecolor{huazul}\specialrule{1.4pt}{0pt}{0pt}
\rowcolor{hufondo}\textbf{Tipo} & \textbf{Cómo responde el estudiante} & \textbf{Cómo se corrige} \\
\arrayrulecolor{hulinea}\midrule
Opción múltiple & Elige una o varias opciones, de texto o de imagen & Una correcta, o varias: todas o al menos una \\
Respuesta corta & Escribe un número o una palabra & Igualdad de texto, sin distinguir mayúsculas ni espacios al inicio o al final \\
Arrastrar y soltar & Lleva cada pieza a su lugar sobre un dibujo & Cada pieza en su lugar; admite otras soluciones y trata como iguales las piezas con la misma imagen \\
Zonas sobre una imagen & Toca una zona del dibujo & La zona elegida está entre las aceptadas \\
Estados por casilla & Elige el estado de cada casilla de una cuadrícula & La cuadrícula completa coincide con una respuesta aceptada \\
Huecos en el texto & Completa cada hueco con una palabra de un banco & El texto completo coincide con una respuesta aceptada \\
\arrayrulecolor{hulinea}\bottomrule
\end{tabularx}
\arrayrulecolor{black}
\caption{Tipos de respuesta y su corrección}
\label{tab:tipos_respuesta}
\end{table}

La dificultad se fija por categoría, porque una misma tarea puede ser difícil para los menores y fácil para los mayores. Las categorías siguen las de Bebras, con nombres de animales de Bolivia, y se definen por cursos: cada una reúne dos cursos seguidos y el estudiante rinde en la de su curso (Tabla~\ref{tab:categorias}).

\begin{table}[H]
\estilotabla
\begin{tabularx}{\linewidth}{l X}
\arrayrulecolor{huazul}\specialrule{1.4pt}{0pt}{0pt}
\rowcolor{hufondo}\textbf{Categoría} & \textbf{Cursos} \\
\arrayrulecolor{hulinea}\midrule
Guacamayo & 1.º y 2.º de primaria \\
Capibara & 3.º y 4.º de primaria \\
Titi & 5.º y 6.º de primaria \\
Jucumari & 1.º y 2.º de secundaria \\
Yaguareté & 3.º y 4.º de secundaria \\
Kuntur & 5.º y 6.º de secundaria \\
\arrayrulecolor{hulinea}\bottomrule
\end{tabularx}
\arrayrulecolor{black}
\caption{Categorías del Desafío Bebras Bolivia}
\label{tab:categorias}
\end{table}

\subsection{Codificación}

El editor de tareas es una sola pantalla para crear y para editar. Arriba se elige cómo responde el estudiante y, según el tipo, cambia solo la sección de la respuesta; el enunciado y la explicación usan el mismo editor de bloques, con texto con formato, imágenes que se agrandan tirando de su borde e imágenes dentro de la línea. Antes de guardar, el editor revisa la tarea y dice qué falta; el servidor repite la misma validación, así que una tarea inválida nunca se guarda.

El probador muestra la tarea con los mismos componentes que ve el estudiante y la corrige en el servidor, con las mismas reglas que se aplicarán en el desafío. Desde el editor, «Probar» lleva lo que está en pantalla sin guardarlo. Las 43 tareas del cuadernillo se cargaron con un validador que comprueba que estén todas, en orden y con sus metadatos; los recortes de sus imágenes se generaron desde el PDF original.

\subsection{Pruebas}

Cada historia se verificó con pruebas automatizadas de extremo a extremo, que recorren las pantallas en el navegador o llaman a la API del servidor, y la lógica de corrección, con pruebas unitarias. La Tabla~\ref{tab:pruebas_i3} resume los casos por tema; el capítulo 9 describe la estrategia completa.

\begin{table}[H]
\estilotabla
\begin{tabularx}{\linewidth}{>{\raggedright\arraybackslash}p{3.6cm} >{\raggedright\arraybackslash}X c c}
\arrayrulecolor{huazul}\specialrule{1.4pt}{0pt}{0pt}
\rowcolor{hufondo}\textbf{Tema} & \textbf{Qué se comprueba} & \textbf{Casos} & \textbf{Resultado} \\
\arrayrulecolor{hulinea}\midrule
Editor y probador & Crear, validar, probar sin guardar, reordenar opciones, filtros del banco y editor en celular & 8 & Aprobado \\
Arrastrar y soltar & Armado sobre el dibujo, gesto en computadora y celular, soluciones alternativas y piezas equivalentes & 14 & Aprobado \\
Zonas sobre una imagen & Dibujo de zonas, deshacer, selección con mouse, teclado y toque, y puntaje en un desafío & 5 & Aprobado \\
Estados y huecos & Configuración, respuesta con banco de palabras, respuestas parciales y las cinco tareas del cuadernillo en desafíos reales & 11 & Aprobado \\
Privacidad y visibilidad & El estudiante nunca recibe la clave; cada nivel de visibilidad decide quién ve la tarea & 5 & Aprobado \\
\arrayrulecolor{hulinea}\bottomrule
\end{tabularx}
\arrayrulecolor{black}
\caption[Pruebas de la Iteración 3]{Pruebas de extremo a extremo de la Iteración 3}
\label{tab:pruebas_i3}
\end{table}

\subsection{Resultados}

La iteración dejó el banco de tareas completo: las 43 tareas del cuadernillo Bebras 2024 están cargadas y cada una se responde como pide su enunciado. El administrador ve todas las tareas en una lista que filtra por tipo de respuesta y por categoría, y edita cada una en una sola pantalla.

La Figura~\ref{fig:desafio_tareas} muestra el trabajo del administrador sobre el banco. En (a) está la lista de las 43 tareas, con los filtros por tipo de respuesta y por categoría; en (b), el editor de una tarea de opción múltiple, donde se redactan el enunciado, las opciones y la explicación.

\parcapturas{assets/ui/desafio/tareas-lista.jpg}{Banco de tareas con sus filtros}{assets/ui/desafio/tarea-opcion-multiple.jpg}{Editor de una tarea de opción múltiple}{Banco de tareas y editor}{fig:desafio_tareas}

El tipo de respuesta cambia solo la sección Respuesta del editor. En arrastrar y soltar, cada pieza se suelta sobre el dibujo en el lugar que le corresponde; en las zonas, se marcan sobre la imagen y la correcta queda en verde. La Figura~\ref{fig:desafio_editores} los muestra: (a) es el editor de arrastrar y soltar y (b), el de zonas sobre una imagen.

\parcapturas{assets/ui/desafio/tarea-arrastrar.jpg}{Arrastrar y soltar}{assets/ui/desafio/tarea-zonas.jpg}{Zonas sobre una imagen}{Editores de respuesta interactiva}{fig:desafio_editores}

En los huecos, el autor selecciona palabras de la pregunta y las convierte en el hueco. Antes de usar una tarea, el probador la muestra tal como la verá el estudiante y la corrige. La Figura~\ref{fig:desafio_probador} presenta en (a) el editor de huecos y en (b) el probador con una tarea ya corregida.

\parcapturas{assets/ui/desafio/tarea-huecos.jpg}{Completar huecos}{assets/ui/desafio/tarea-probador.jpg}{Probador con la tarea corregida}{Huecos en el texto y probador}{fig:desafio_probador}

\section{ITERACIÓN 4}

La Iteración 4 construyó los desafíos con los que se usa ese banco. Se describen su planificación, el diseño del ciclo de un desafío, la codificación, las pruebas y los resultados.

\subsection{Planificación de la iteración}

Con el banco listo, la iteración organizó los desafíos: un mismo desafío reúne varias categorías, cada una con sus preguntas, y avanza solo de la inscripción a la rendición según sus fechas. Se sumaron la práctica pública y la portada, que muestra al público los desafíos vigentes.

La Tabla~\ref{tab:hus_iteracion_4} lista las historias de la iteración; sus tarjetas, con los criterios de aceptación, están en el Anexo A.

\begin{table}[H]
\estilotabla
\begin{tabularx}{\linewidth}{l >{\raggedright\arraybackslash}X >{\raggedright\arraybackslash}p{2.8cm} c c}
\arrayrulecolor{huazul}\specialrule{1.4pt}{0pt}{0pt}
\rowcolor{hufondo}\textbf{ID} & \textbf{Historia} & \textbf{Usuario} & \textbf{Prioridad} & \textbf{Esfuerzo} \\
\arrayrulecolor{hulinea}\midrule
HU-31 & Crear un desafío & Administrador & 1 & 5 \\
HU-32 & Elegir las preguntas de cada categoría & Administrador & 1 & 8 \\
HU-33 & Programar las fechas del desafío & Administrador & 1 & 8 \\
HU-34 & Ajustar el puntaje & Administrador & 2 & 3 \\
HU-35 & Decidir qué ve el estudiante al terminar & Administrador & 2 & 5 \\
HU-36 & Mezclar el orden de las opciones & Administrador & 2 & 3 \\
HU-37 & Publicar el desafío & Administrador & 1 & 5 \\
HU-38 & Recorrer el desafío como un estudiante & Administrador & 2 & 5 \\
HU-39 & Pausar un desafío en curso & Administrador & 2 & 8 \\
HU-40 & Practicar sin crear una cuenta & Estudiante & 1 & 8 \\
HU-41 & Consultar el avance de la práctica & Estudiante & 2 & 3 \\
HU-42 & Armar una práctica para un grupo & Maestro & 2 & 8 \\
HU-43 & Rendir la práctica de un maestro & Estudiante & 2 & 3 \\
HU-44 & Ver qué desafíos hay & Visitante & 1 & 5 \\
HU-45 & Ir directo a la acción de cada desafío & Visitante & 2 & 3 \\
\arrayrulecolor{hulinea}\midrule
\rowcolor{white}\multicolumn{4}{r}{\textbf{Total de puntos}} & \textbf{80} \\
\arrayrulecolor{hulinea}\bottomrule
\end{tabularx}
\arrayrulecolor{black}
\caption[Historias de la Iteración 4]{Historias de la Iteración 4}
\label{tab:hus_iteracion_4}
\end{table}


\subsection{Diseño}

El estado de un desafío no se guarda: se calcula con sus fechas y con las acciones del administrador. Así el desafío avanza solo, sin procesos programados que puedan fallar. La Figura~\ref{fig:estados_desafio} muestra sus estados como cajas y los cambios entre ellos como flechas rotuladas con la acción que los provoca. Las flechas naranjas son acciones del administrador; las punteadas ocurren solas, según las fechas. Se lee desde «Borrador»: al publicar pasa a «Inscripción abierta» o, si tiene fecha de apertura, a «Programado». Luego sigue la fila superior hasta «En curso», que puede pausarse y reanudarse, y la fila inferior de derecha a izquierda hasta «Resultados publicados».

\begin{figure}[H]
\centering
\begin{tikzpicture}[ancho/.style={minimum width=2.4cm, text width=2.1cm}]
\node[inicio] (ini) at (0,1.5) {};
\node[estespera, ancho] (borrador) at (0,0) {Borrador};
\node[estespera, ancho] (programado) at (0,-2.3) {Programado};
\node[estactivo, ancho] (inscripcion) at (4,0) {Inscripción abierta};
\node[estactivo, ancho] (preparacion) at (8,0) {Inscripción cerrada};
\node[estactivo, ancho] (curso) at (12,0) {En curso};
\node[estalerta, ancho] (pausado) at (12,-2.3) {Pausado};
\node[estfinal, ancho] (terminado) at (12,-4.6) {Terminado};
\node[estfinal, ancho] (listos) at (8,-4.6) {Resultados listos};
\node[estfinal, ancho, line width=1.6pt] (publicados) at (4,-4.6) {Resultados publicados};
\draw[flecha] (ini) -- node[etq, right=3pt]{crear} (borrador);
\draw[flechaadm] (borrador) -- node[etq, above=2pt]{publicar} (inscripcion);
\draw[flechaadm] (borrador) -- node[etq, left=3pt]{publicar} (programado);
\draw[flechaauto] (programado.north east) -- node[etq, sloped, below=2pt]{apertura} (inscripcion.south west);
\draw[flechaauto] (inscripcion) -- node[etq, above=2pt]{cierre} (preparacion);
\draw[flechaauto] (preparacion) -- node[etq, above=2pt]{inicio} (curso);
\draw[flechaadm] ([xshift=-5pt]curso.south) -- node[etq, left=3pt]{pausar} ([xshift=-5pt]pausado.north);
\draw[flechaadm] ([xshift=5pt]pausado.north) -- node[etq, right=3pt]{reanudar} ([xshift=5pt]curso.south);
\draw[flechaauto] (curso.east) -- ++(0.65,0) |- node[etq, pos=0.25, right=3pt]{cierre} (terminado.east);
\draw[flechaauto] (terminado) -- node[etq, above=2pt]{calcular} (listos);
\draw[flechaauto] (listos) -- node[etq, above=2pt]{publicar} (publicados);
\end{tikzpicture}
\caption[Estados de un desafío]{Estados de un desafío}
\label{fig:estados_desafio}
\end{figure}

El paso del tiempo mueve el desafío entre inscripción, rendición y cierre; el administrador lo publica, lo pausa y lo reanuda. Al publicarse queda programado si su inscripción todavía no empieza. Un desafío en curso o terminado no se puede editar, y solo se elimina en borrador. El cálculo y la publicación de los resultados ocurren solos al cierre o en la fecha de resultados, y el administrador también puede hacerlos a mano; se describen en el módulo de evaluación (capítulo 8).

El puntaje sigue el estándar de Bebras y se puede ajustar en cada desafío. Cada categoría parte de un puntaje inicial igual a la suma de todas las penalizaciones posibles, así el total nunca baja de cero.

\begin{table}[H]
\estilotabla
\begin{tabularx}{\linewidth}{X c c c}
\arrayrulecolor{huazul}\specialrule{1.4pt}{0pt}{0pt}
\rowcolor{hufondo}\textbf{Dificultad} & \textbf{Correcta} & \textbf{Incorrecta} & \textbf{Sin responder} \\
\arrayrulecolor{hulinea}\midrule
Fácil & +6 & $-2$ & 0 \\
Medio & +9 & $-3$ & 0 \\
Difícil & +12 & $-4$ & 0 \\
\arrayrulecolor{hulinea}\bottomrule
\end{tabularx}
\arrayrulecolor{black}
\caption{Puntaje estándar por dificultad}
\label{tab:puntaje}
\end{table}

\subsection{Codificación}

El desafío se edita en una sola pantalla, con sus categorías, sus fechas, su puntaje y lo que verá el estudiante; las preguntas de cada categoría se eligen en otra, con los mismos filtros del banco. Al agregar una pregunta se copian su dificultad y su puntaje; el enunciado y la respuesta siguen siendo los de la tarea del banco, que no se puede editar mientras el desafío está en curso o pausado. Las reglas del desafío (estados, ventanas, categorías, cursos y puntaje) viven en un módulo propio que recibe la hora como dato, lo que permite probarlas sin servidor.

La práctica pública lista las tareas liberadas de cada categoría, de la más fácil a la más difícil, las corrige en el servidor y recuerda cuáles resolvió el estudiante, que no necesita cuenta, solo mientras la pestaña del navegador sigue abierta: al cerrarla, el avance se borra y no queda en un equipo compartido. El maestro puede armar además una práctica con tareas liberadas para sus estudiantes.

\subsection{Pruebas}

Como en la iteración anterior, cada historia se verificó con pruebas de extremo a extremo en el navegador o contra la API, y las reglas del desafío, con pruebas unitarias que fijan la hora. La Tabla~\ref{tab:pruebas_i4} agrupa los casos por tema.

\begin{table}[H]
\estilotabla
\begin{tabularx}{\linewidth}{>{\raggedright\arraybackslash}p{3.6cm} >{\raggedright\arraybackslash}X c c}
\arrayrulecolor{huazul}\specialrule{1.4pt}{0pt}{0pt}
\rowcolor{hufondo}\textbf{Tema} & \textbf{Qué se comprueba} & \textbf{Casos} & \textbf{Resultado} \\
\arrayrulecolor{hulinea}\midrule
Categorías y fechas & Un desafío con varias categorías, cursos válidos de cada una y fases programadas & 5 & Aprobado \\
Protección y pausa & Un desafío en curso no se edita, la pausa bloquea a los estudiantes y devuelve el tiempo, y no se borra lo ya rendido & 8 & Aprobado \\
Puntaje & Escalas de Bebras y la distribución estándar de 0 a 180 puntos & 2 & Aprobado \\
Probar como estudiante & Recorrer el desafío sin dejar rastro & 1 & Aprobado \\
Práctica & Corrección de todos los tipos, accesibilidad, orden por dificultad, avance guardado y prácticas del maestro & 15 & Aprobado \\
Portada y escala & Solo desafíos publicados, acciones según la etapa y listados con más de cien registros & 3 & Aprobado \\
\arrayrulecolor{hulinea}\bottomrule
\end{tabularx}
\arrayrulecolor{black}
\caption[Pruebas de la Iteración 4]{Pruebas de extremo a extremo de la Iteración 4}
\label{tab:pruebas_i4}
\end{table}

\subsection{Resultados}

Un desafío reúne varias categorías, cada una con sus preguntas, y su estado (borrador, inscripción abierta, en curso o resultados publicados) se ve en la lista.

La Figura~\ref{fig:desafio_gestion} corresponde a la gestión de desafíos. En (a) se ve la lista, con el estado de cada uno, y en (b) el editor, que organiza el desafío por categorías.

\parcapturas{assets/ui/desafio/desafios-lista.jpg}{Lista de desafíos}{assets/ui/desafio/desafio-editor.jpg}{Editor del desafío por categorías}{Gestión de desafíos}{fig:desafio_gestion}

Las preguntas de cada categoría se eligen del banco con los mismos filtros, y el sistema avisa si una pregunta ya es pública en la práctica. La portada muestra al público los desafíos vigentes con lo que cada visitante puede hacer en ese momento.

En la Figura~\ref{fig:desafio_portada}, (a) es la elección de las preguntas de una categoría desde el banco y (b), la portada pública con los desafíos vigentes.

\parcapturas{assets/ui/desafio/desafio-preguntas.jpg}{Elección de preguntas}{assets/ui/desafio/portada.jpg}{Portada con los desafíos}{Preguntas del desafío y portada}{fig:desafio_portada}
